\subsubsection{Brookings: Multi-Dimensional AI Competition and China's Multi-Track Strategy}

Two Brookings reports published in April 2026 provide important support for our analytical framework. \textit{Competing AI Strategies for the US and China} argues that the US-China AI race has evolved from a single-dimensional contest of model performance into a five-dimensional competition spanning compute, model capability, industry adoption, system integration, and global deployment. In this multi-dimensional framework, the US and China hold advantages in \textit{different dimensions}: the US leads in frontier model capability, while China leads substantially in application diffusion speed and open-source ecosystem penetration.

More insightful is the core thesis of \textit{China is running multiple AI races}: Chinese AI developers are not simply chasing GPT-5---they are running \textit{different races} altogether. Driven by both industrial constraints (chip export controls) and Beijing's policy focus (AI empowering the real economy), China's AI industry has built a differentiated competitive edge centered on open-source ecosystems, industrial applications, and cost advantages. This aligns perfectly with our spectrum diffusion framework: lowering access barriers and accelerating spectrum ascent has greater societal impact than pursuing the extreme performance of any single model.

Brookings Fellow Kyle Chan confirmed this trend in a July 2026 CNBC interview: `Chinese AI models are particularly attractive to American companies now as AI costs skyrocket. Where previously U.S. companies were prioritizing AI adoption regardless of model, now they're getting more cost-conscious.'' He estimated that Chinese frontier open-source models currently lag top US models by `six to nine months''---a gap that is rapidly narrowing in the face of cost advantages.

\subsubsection{RAND: Open Source as Soft Power and the AI Competition Spectrum}

RAND's \textit{Open Models, Soft Power, and the Spectrum of U.S.-China Artificial Intelligence Competition} presents an analytical framework highly resonant with ours. The report makes two core contributions.

First, it positions open-source AI as a \textbf{soft power instrument}. China, through open-source models (DeepSeek, GLM, Qwen, etc.), expands its influence over Global South nations that cannot afford OpenAI's or Anthropic's high API fees but can gain AI capabilities by deploying Chinese open-source models. This is cognitive democratization manifesting at the international system level: open source transforms AI capability from a privilege of a few wealthy nations into a global public good.

Second, the US-China AI competition is not a binary opposition but a \textbf{spectrum}---a continuous distribution from fully open to fully closed. China sits at the open end of the spectrum; the US moved toward the closed end under Biden (driven by safety executive orders) and swung back toward the open end under Trump's second term. This very policy oscillation---RAND notes---may damage US credibility and leadership in the open-source ecosystem.

RAND's policy recommendation carries direct relevance for our argument: \textbf{if the US excessively restricts open source, it risks ceding the global AI ecosystem to China.}

\subsubsection{Goldman Sachs AI Adoption Tracker: Empirical Anchors for Spectrum Distribution}

Goldman Sachs's AI Adoption Tracker, released in April 2026 and based on the US Census Bureau's Business Trends and Outlook Survey (BTOS), provides the closest thing to systematic empirical data supporting our spectrum distribution hypothesis:

\begin{center}
\begin{tabular}{c|c}
\hline
Metric & Data \\
\hline
Overall US firm AI adoption rate & 19.8\% (projected 22.3\% in 6 months) \\
Large firms ($> employees) adoption & 35.3\% \\
Small firms (20--49 employees) recent growth & +2.1 pp to 21.5\% \\
IT/Web services sector adoption & 60\% (highest) \\
Daily time saved per AI-using worker & 40--60 minutes (OpenAI enterprise data) \\
Academic studies: productivity uplift & 23\% average \\
Company anecdotes: efficiency gains & $\sim\% \\
\hline
\end{tabular}
\end{center}

These data provide threefold validation for our theoretical framework:

\begin{enumerate}
\item \textbf{Non-uniform spectrum distribution receives preliminary empirical support}: only 19.8\% of US firms use AI---meaning over 80\% remain at what we define as Level 0 (no use at all). Among those that do, the gap between 60\% adoption (IT sector) and the lowest sectoral rates corresponds precisely to the differentiation from Levels 1--2 to Levels 4--5 in our spectrum.
\item \textbf{Firm size correlates positively with spectrum level}: the 35.3\% vs.\ 21.5\% adoption gap between large and small firms supports our claim that AI use levels are highly correlated with existing social stratification. Large firms have more resources for AI training, deployment, and customization---the external driving forces needed for spectrum ascent.
\item \textbf{Productivity gains validate the value of spectrum ascent}: 23--33\% efficiency gains and 40--60 minutes saved per day demonstrate that transitioning from Level 0 to Level 2+ yields significant economic returns. This provides a cost-benefit rationale for our policy recommendation of lowering spectrum ascent barriers through education and open source.
\end{enumerate}

Goldman's economists wrote: `We continue to observe large impacts on labor productivity in the limited areas where generative AI has been deployed''---and then identified the core contradiction: `The companies using AI are pulling ahead, and most of their competitors aren't even in the race yet.'' This is the precise industrial-level mapping of the spectrum solidification mechanism described in Section 2.6.

\subsubsection{Chinese Official Data: Industry Scale and Penetration Targets}

The China Academy of Information and Communications Technology (CAICT) February 2026 White Paper on Global AI Development provides macro data from the Chinese side: China's AI core industry exceeded 900 billion yuan in 2024, large model comprehensive performance has significantly improved, and application barriers and usage costs continue to decline. The White Paper identifies open source as a critical pathway for global AI development and notes that China ranks second globally in the number of open-source large models.

More strategically significant is the State Council's August 2025 `Opinions on Deepening the Implementation of the AI+ Action Initiative,'' which set explicit quantitative targets:
\begin{itemize}
\item \textbf{By 2027}: AI deeply integrated across 6 key sectors; next-generation intelligent terminals and AI agents surpass 70\% adoption
\item \textbf{By 2030}: AI comprehensively empowers high-quality development nationwide
\end{itemize}

A 70\% penetration target---if achieved---would mean that under policy-driven impetus, Chinese society's AI cognitive spectrum evolves rapidly from the current pyramid shape (majority at Levels 0--2) toward a more balanced distribution. This is the national-policy-level manifestation of our `optimized spectrum diffusion'' scenario (Scenario B). To be clear, policy targets are not reality---but policy targets shape the direction of resource allocation, and the direction of resource allocation affects the trajectory of spectrum evolution.

\subsection{Synthesis: Multi-Source Empirical Support for the Theoretical Framework}

Mapping the latest think-tank and financial-institution findings onto our theoretical framework yields a multi-source validation matrix:

\begin{center}
\begin{tabular}{c|c|c}
\hline
Core Thesis & Validation Source & Type \\
\hline
Non-uniform spectrum distribution & Goldman Sachs 19.8\% adoption data & Quantitative empirical \\
Open source lowers spectrum ascent barriers & CNBC: Chinese models 60--90\% cheaper; US token share surged 4.5\%$\rightarrow\% & Market behavior \\
Closed source creates gatekeeper structures & RAND: open source as soft power; OpenAI limits releases at government request & Policy analysis \\
Spectrum ascent requires external driving forces & Goldman: large firms 35.3\% vs.\ small 21.5\% & Structural \\
Cognitive democratization as international competition & Brookings: multi-dimensional AI race; State Council: 70\% target & Strategic analysis \\
Geotechnological significance of open source & RAND: US risks ceding ecosystem to China if restricting open source & Geopolitical \\
\hline
\end{tabular}
\end{center}

The significance of this multi-source validation matrix is that our core argument receives independent support across four dimensions: quantitative empirical evidence (Goldman Sachs data), market behavior (CNBC reporting), policy analysis (Brookings/RAND), and strategic planning (Chinese State Council targets). This not only strengthens the credibility of our argument but also demonstrates that our theoretical framework possesses analytical applicability across micro (firm adoption rates), meso (industrial competition), and macro (national strategy) levels.
